\section{How to use: first example}

You need to load the \NamePack{alterqcm.sty} with  |\usepackage[english]{alterqcm}|, if you want to use the english language. With some languages like Greek or Chinese you need to compile with Xe\LaTeX\ otherwise you can compile with Lua\LaTeX\ or PDF\LaTeX\ .

Just use an environment \tkzname{alterqcm} and the macro \tkzcname{AQquestion}, here is an example :


 \noindent
\begin{minipage}[c][][t]{.40\linewidth}
\begin{tkzexample}[code only,small]
 \documentclass[12pt]{article}
 \usepackage[english]{alterqcm}
 % or french ...
 \usepackage{fullpage}
 \parindent0pt
 \begin{document}
 \begin{alterqcm}
  \AQquestion{Question}{%
  {Proposition 1},
  {Proposition 2},
  {Proposition 3}}
 \end{alterqcm}
 \end{document}\end{tkzexample}
\end{minipage}\hfill \noindent
\begin{minipage}[c][][b]{.50\linewidth}
\textbf{alterqcm.sty} creates a new environment \textbf{alterqcm} which allows for a two-column table. One column on the left for the questions, the other for the different proposals.  The propositions are given in a list :

\tkzname{\{\{Proposition 1\},\\\{Proposition 2\},\\\{Proposition 3\}\}}.

 The number of propositions is between \tkzname{2} and \tkzname{5}.
\end{minipage}

\medskip
The result is:

\bigskip
  \begin{alterqcm}
  \AQquestion{Question}
  {%
  {Proposition 1},
  {Proposition 2},
  {Proposition 3}%
  }
  \end{alterqcm}

\medskip
 The total width of the array is equal to \tkzcname{textwidth}. By default the question column has the width \tkzname{100mm} plus a few millimeters ... introduced by the table. The width of the answers is equal to \tkzcname{textwidth} minus the width of the first column. \Imacro{textwidth}

The important point is that the height of the lines in the proposals is calculated automatically so that, on the one hand, the text of the proposals is placed correctly without touching the lines and, on the other hand, the text of the corresponding question can be included in its box. Precise positioning is obtained with the option \tkzname{pq}.

\subsection{Packages loaded by \tkzname{alterqcm.sty}}
The list of loaded packages is as follows:

\begin{tkzexample}[code only]
  \RequirePackage{xkeyval}[2005/11/25]
  \RequirePackage{calc}
  \RequirePackage{ifthen,forloop}
  \RequirePackage{array}
  \RequirePackage{multirow}
  \RequirePackage{pifont}
\end{tkzexample}


\tkzHandBomb\ You will need to load \tkzname{longtable.sty} if you wish to use the \tkzname{long} option for one of your arrays.

\tkzHandBomb\ You also need the macro \tkzcname{square}, it is either defined in the package \tkzname{fourier} or in the package.
\endinput